% Generated draft: requires mathematical and source review.
\documentclass[11pt]{article}
\usepackage[margin=25mm]{geometry}
\usepackage{fontspec}
\setmainfont{Latin Modern Roman}
\usepackage{amsmath,amssymb,mathtools,unicode-math}
\setmathfont{Latin Modern Math}
\usepackage{enumitem,needspace,xurl,hyperref}
\setlength{\parindent}{0pt}
\setlength{\parskip}{5pt}
\setlength{\emergencystretch}{3em}
\newcommand{\N}{\mathbb{N}}
\newcommand{\Z}{\mathbb{Z}}
\newcommand{\Q}{\mathbb{Q}}
\newcommand{\R}{\mathbb{R}}
\newcommand{\C}{\mathbb{C}}
\newcommand{\F}{\mathbb{F}}
\newcommand{\A}{\mathbb{A}}
\newcommand{\PP}{\mathbb{P}}
\newcommand{\E}{\mathbb{E}}
\newcommand{\eps}{\varepsilon}
\newcommand{\dd}{\,\mathrm{d}}
\begin{document}
\section*{UnsolvedMath ID: 30000462}
\textbf{Source problem code:} catalogue entry
\subsection{Definitions and mathematical statement}
\textbf{Answered example, not an open conjecture.} The supplied verification marks this extraction unrecoverable as an open problem: the official report presents an elementary street-grid example and immediately gives the answer
\[
 \binom{a+b}{a}-\binom{a+b}{a-1}.
\]
The retained original instead mentions a ``relevant subspace lattice'' without defining it. No subspace-lattice or finite-field analogue is inferred from the grid answer.

The standard grid formulation consistent with that answer is as follows. For integers $0\leq a\leq b$, count shortest routes from $(0,0)$ to $(a,b)$ with steps $(1,0)$ and $(0,1)$ such that each visited vertex $(i,j)$ satisfies $j\geq i$. Every such route has length $a+b$. In this explicit interpretation the stated answer is
\[
 N(a,b)=\binom{a+b}{a}-\binom{a+b}{a-1},
\]
with the convention $\binom Nk=0$ when $k<0$ or $k>N$. The endpoint restriction $a\leq b$ is essential for this formula as a nonnegative count; if $a>b$, no route satisfies the terminal inequality. The excerpt itself does not state the grid endpoints, step set or parameter range, so this is a qualified reconstruction of the \emph{answered} example rather than a new solved theorem about an unspecified subspace lattice.

The literature triage checked 21 August 2026 says the grid and boundary definitions are missing and no resolution was identified, whereas the verification explicitly reports the answer in the source. Both pieces of metadata are retained; the registry open label is not promoted over the explicit answered-example qualification.
\subsection{Short English statement}
This is a street-grid counting example already answered in the source, not an independent open problem. Under the usual route interpretation its count is $\binom{a+b}{a}-\binom{a+b}{a-1}$ for $0\leq a\leq b$. The original's undefined subspace-lattice wording cannot be substituted for that model.
\subsection{Sources}
\begin{itemize}
\item Local UnsolvedMath registry, permanent ID 30000462, exact source code catalogue entry; \emph{Answered Shortest-Route Enumeration}. Proposed year 2006, registry status open, statement\_status unrecoverable. Supplied verification says the official report immediately gives the displayed binomial answer and that this is not an open problem; its claimed inspection is metadata.
\item Supplied report citation: \emph{Representations of Finite Groups (2007)}, Oberwolfach Reports, DOI 10.4171/owr/2006/15, pp. 895--978, \url{https://doi.org/10.4171/owr/2006/15}. Title year differs from DOI/proposed year.
\item Supplied literature triage checked 2026-08-21, citing only ``original Oberwolfach report'', \url{https://doi.org/10.4171/owr/2006/15}. Its missing-grid/no-resolution assessment is retained alongside the verification's answered-example correction. The dated OWR-derived curation requests independent verification of time-sensitive status and cited claims.
\end{itemize}
\end{document}
